\clearpage
\fancyhead[L]{\small PASQUICH. EHRENRETTUNG.}
\hypertarget{native-639-lower}{}
\expandafter\def\csname HistoricalNoteText35\endcsname{Der Collimationsfehler des Königsberger Meridiankreises beträgt beinahe 2 Grad. Gewöhnlich bringen die Künstler bei den Instrumenten Vorrichtungen an, die Collimationsfehler wegzuschaffen: der Astronom weiss es ihnen keinen Dank. Dass an einigen Meridiankreisen das Fadennetz auch im verticalen Sinn beweglich ist, hat man nicht wie eine Vollkommenheit anzusehen.\hfill G.}
\expandafter\def\csname HistoricalNoteText36\endcsname{Bei der ersten Beobachtung habe ich mich an den Stundenwinkel gehalten, da die Rectascension offenbar durch einen Druck- oder Schreibfehler entstellt ist, und anstatt $357^\circ\,49'\,14''5$ sein sollte $357^\circ\,49'\,4''5$.}
\begin{center}\small Astronomische Nachrichten. Band III. Nr. 53. Seite 77..84. 1824 März 3.\\[3pt]\rule{.72\linewidth}{.3pt}\end{center}
\begin{center}\emph{Ehrenrettung.}\end{center}

In einem der letzten Stücke der in Genua erscheinenden Correspondance Astronomique ist der Professor \textsc{Pasquich} angeklagt, diejenigen Positionen des Cometen von 1821, welche in Nr. 2 der Astr. Nachr. abgedruckt sind, nicht aus Beobachtungen am Aequatoreal, wie \textsc{Pasquich} selbst versichert hat, abgeleitet, sondern sie erdichtet, d. i. aus Elementen berechnet zu haben. Zur Begründung dieser Anklage wird nicht etwa behauptet, dass an den Tagen, für welche die Cometen-Positionen bekannt gemacht waren, auf der Ofener Sternwarte gar nicht, oder dass unrichtig beobachtet wäre, sondern im Gegentheil, der Ankläger ist selbst ein Theilnehmer an den Beobachtungen gewesen, und er bringt selbst einen Theil der Originalbeobachtungen aus \textsc{Pasquich}'s eigner vom Notarius beglaubigten Handschrift bei. Auch beweiset und, genau genommen, behauptet der Ankläger nicht einmal, dass aus den Originalbeobachtungen, wenn sie richtig und vollständig reducirt würden, andere Positionen folgten, als die von \textsc{Pasquich} bekannt gemachten. Worauf ist denn also, wird man fragen, eine Anklage gegründet, wodurch der Untergebene seinen Vorgesetzten als eines Verbrechens schuldig darstellt, welches das entehrendste ist, das einem Astronomen vorgeworfen werden kann? Lediglich darauf, dass das Instrument gar nicht berichtigt gewesen sei, dass aus Beobachtungen an einem nicht berichtigten Instrumente keine richtigen Positionen abgeleitet werden können, und dass also die von \textsc{Pasquich} bekannt gemachten Positionen, deren Richtigkeit sofort durch die Rechnung des Dr. \textsc{Ursin} bestätigt war, auf eine unredliche Art fingirt sein müssten.

Es bedarf für Astronomen blos dieser einfachen Darstellung, oder einer nur einigermassen aufmerksamen Durchlesung der Anklage selbst, um einzusehen, dass diese durchaus aller Begründung ermangelt. Jeder Astronom weiss, dass aus Beobachtungen an einem nicht berichtigten Instrumente ebenso zuverlässige Positionen abgeleitet werden können, wie aus Beobachtungen an einem vollkommen berichtigten. Man bestimmt die Grösse der Abweichungen des Instruments vom vollkommen berichtigten Zustande durch schickliche Beobachtungen und bringt die daraus folgenden Correctionen bei allen andern Beobachtungen in Rechnung. Der Calcul ist allemal die schärfste Berichtigungsart. Ob die \SourceLeaf{640} Correctionen einige Secunden oder einige Minuten betragen, ist für die Genauigkeit der Resultate gleichgültig. Nicht um dieser willen, sondern der Bequemlichkeit wegen, zieht man im Allgemeinen vor, nur mit kleinen Berichtigungen zu thun zu haben, und bei den immer gleich viel wirkenden 'Abweichungen ist auch hieran so viel wie gar nichts gelegen\HistoricalNote{35}.

Die Behauptung des Anklägers, man könne aus Beobachtungen an einem nicht berichtigten Instrumente keine richtige Positionen ableiten, läuft demnach nur auf ein Geständniss hinaus, welches keines Commentars bedarf. Seinem Vorgesetzten, der seit langer Zeit als Mathematiker vortheilhaft bekannt ist, konnte natürlich eine Correctionsrechnung nicht schwer fallen, die ja im Bereich eines nur einigermassen geübten Anfängers ist.

Wenn man die Anklage nur oberflächlich liest, könnte man vielleicht glauben, dass das Instrument sehr weit von dem vollkommen berichtigten Zustande entfernt gewesen sei. Allein

\emph{Erstens} folgt dies nicht aus den angeblichen Beweisen. Es wird erzählt, dass einst in der Dämmerung \textsc{Pasquich} den Cometen mit dem Aequatoreal nicht finden konnte, als der Ankläger ihn schon im Cometensucher sah. Dies beweist gar nichts. Der Comet konnte wirklich im Felde, und doch, eben wegen der Dämmerung, mit der starken Vergrösserung noch nicht erkennbar sein, als ein lichtstarker Cometensucher mit schwächerer Vergrösserung ihn schon zeigte; eben so beweisen des Anklägers Rechnungen, die er über die Beobachtungen geführt hat, an sich noch gar nichts für eine mangelhafte Berichtigung: in der That hätte er mit \emph{solchen} Rechnungen eine Abweichung finden müssen, wenn gar keine vorhanden war. Er vernachlässigt nemlich ganz die Refraction und macht sich, in seiner Art, lustig über \textsc{Pasquich}'s Versicherung, diese bei der Reduction der Beobachtungen berücksichtigt zu haben, die nach des Anklägers Behauptung im Maximum weder bei der Differenz der Rectascensionen, noch bei der der Declinationen auf zwei Secunden steigen könne. Wusste der Ankläger auch diese leichte Rechnung nicht zu führen? Die Differenz beträgt am 27sten Februar bei beiden über eine Minute.

\emph{Zweitens} aber lässt sich wirklich aus den aufgestellten Beobachtungen hinreichend erkennen, dass die Abweichung des Instruments keinesweges so enorm gross gewesen ist.

Die wenigen Beobachtungen von $\gamma$ Pegasi, welche der Ankläger zu unserer Kenntniss gebracht hat, sind zwar unzulänglich zu einer scharfen und vollständigen Bestimmung der Corrections-Elemente des Instruments; sie reichen aber hin, um zu beweisen, dass die Abweichungen nicht so gross sind, um das Auffinden von Sternen zu erschweren; sie reichen ferner hin, um die Cometenbeobachtungen selbst sehr nahe zu reduciren. Ich theile die Resultate meiner darüber geführten Rechnung mit desto grösserem Vergnügen mit, da daraus mit aller nur zu wünschenden Evidenz hervorgeht, dass die Beschuldigung mehr als grundlos, dass sie falsch ist.'

Die Discussion aller 5 Beobachtungen von $\gamma$ Pegasi gab mir
\begin{center}\footnotesize\setlength{\tabcolsep}{3pt}
\begin{tabular}{@{}l r@{}}
Entfernung des Pols des Instruments vom wahren Weltpol & $=5'\;5''56$\\
Stundenwinkel des erstern . . . . . . . . . . . & $212^\circ\;27'\;27''$\\
Correction des Index für den Stundenwinkel . . . . . & $3'\;41''57$\\
für die Declination . . . . . . . . . . . . . . & $4\;57.55$
\end{tabular}\end{center}

Ob und wie viel die optische Axe von dem Parallelismus mit dem Declinationskreise, und dieser von der Verticalität zum Aequatorskreise abweiche, lässt sich aus den Beobachtungen Eines Sterns nicht \SourceLeaf{641} bestimmen; auf die Reduction der Cometenörter kann 'dies aber bei der geringen Verschiedenheit der Declinationen keinen bemerkbaren Einfluss haben.

Die 5 Beobachtungen von $\gamma$ Pegasi werden mit diesen Elementen folgendermaassen
dargestellt:

\begin{center}\footnotesize\setlength{\tabcolsep}{4pt}
\begin{tabular}{r|r|r|l|r}
 & \multicolumn{2}{c|}{Corr. d. St. W. wegen} & Corrig. Beob. & Untersch.\\
 & des Instr. & der Refr. & Stundenwinkel & v. wahren \\
\hline
Febr. 20 & $-162''99$ & $+147''26$ & $81^\circ\,36'\,45''01$ & $+5''42$ \\
22 & $-174.81$ & $+95.54$ & $69\ \ 38\ \ 14.47$ & $\mathord{\cdot}\,+3.23$ \\
22 & $-168.76$ & $+116.59$ & $75\ \ 29\ \ 8.57$ & $+0.33$ \\
27 & $-174.97$ & $+95.56$ & $69\ \ 28\ \ 41.31$ & $+1.42$ \\
27 & $-168.59$ & $+117.84$ & $75\ \ 39\ \ 9.99$ & $-10.40$ \\
\end{tabular}
\end{center}

\begin{center}\footnotesize\setlength{\tabcolsep}{4pt}
\begin{tabular}{r|r|r|l|r}
 & \multicolumn{2}{c|}{Corr. d. Decl. wegen} & Corrig. beobach- & Untersch.\\
 & des Instr. & der Refr. & tete Decl. & v. d. wahren \\
\hline
Febr. 20 & $-97''54$ & $-147''48$ & $14^\circ\,11'\,19''98$ & $-0''63$ \\
22 & $-54.04$ & $-96.10$ & $14\ \ 11\ \ 21.86$ & $+1.42$ \\
22 & $-74.06$ & $-116.33$ & $14\ \ 11\ \ 17.61$ & $-2.83$ \\
27 & $-53.50$ & $-96.15$ & $14\ \ 11\ \ 20.35$ & $+0.33$ \\
27 & $-74.68$ & $-117.59$ & $14\ \ 11\ \ 21.73$ & $+1.71$ \\
\end{tabular}
\end{center}

Wendet man dieselben Elemente zur Reduction der Cometenbeobachtungen an, so erhält
man

\begin{center}\footnotesize\setlength{\tabcolsep}{4pt}
\emph{Für den Stundenwinkel}
\end{center}

\begin{center}\footnotesize\setlength{\tabcolsep}{4pt}
\begin{tabular}{r|r|r|l|r}
 & \multicolumn{2}{c|}{Correction wegen} & Corrigirter & \shortstack{Abw. d. Mitt.\\v. \textsc{Pasquich}'s}\\
 & des Instr. & der Refr. & Stundenwinkel & Angaben \\
\hline
Febr. 22 & $-171''15$ & $+101''34$ & $71^\circ\,48'\,50''19$ & \raisebox{-5.5pt}[0pt][0pt]{$-4''34$} \\
22 & $-165.09$ & $+124.21$ & $77\ \ 39\ \ 19.12$ & \\
26 & $-164.37$ & $+134.96$ & $79\ \ 28\ \ 15.59$ & \raisebox{-5.5pt}[0pt][0pt]{$+1.37$} \\
26 & $-161.37$ & $+154.53$ & $82\ \ 44\ \ 53.16$ & \\
27 & $-161.34$ & $+157.80$ & $83\ \ \ 9\ \ 56.48$ & \raisebox{-5.5pt}[0pt][0pt]{$+4.95$} \\
27 & $-157.98$ & $+191.42$ & $87\ \ 10\ \ 33.44$ & \\
\end{tabular}
\end{center}

\begin{center}\footnotesize\setlength{\tabcolsep}{4pt}
\emph{Für die Declination}
\end{center}

\begin{center}\footnotesize\setlength{\tabcolsep}{4pt}
\begin{tabular}{r|r|r|l|r}
 & \multicolumn{2}{c|}{Correction wegen} & Corrigirte & \shortstack{Abw. d. Mitt.\\v. \textsc{Pasquich}'s}\\
 & des Instr. & der Refr. & Declination & Angaben \\
\hline
Febr. 22 & $-61''15$ & $-100''78$ & $14^\circ\,38'\,10''07$ & \raisebox{-5.5pt}[0pt][0pt]{$+7''42$} \\
22 & $-82.05$ & $-123.18$ & $14\ \ 38\ \ 24.77$ & \\
26 & $-89.01$ & $-134.43$ & $14\ \ 23\ \ \ 9.56$ & \raisebox{-5.5pt}[0pt][0pt]{$+0.91$} \\
26 & $-102.09$ & $-154.65$ & $14\ \ 23\ \ 18.26$ & \\
27 & $-103.80$ & $-158.18$ & $14\ \ 18\ \ 44.02$ & \raisebox{-5.5pt}[0pt][0pt]{$-5.58$} \\
27 & $-120.75$ & $-193.94$ & $14\ \ 18\ \ 34.82$ & \\
\end{tabular}
\end{center}

Für jeden, welcher im astronomischen Calcul kein Fremdling ist, müssen diese
Abweichungen der klarste Beweis sein, dass \textsc{Pasquich}'s Positionen wirklich
aus den Beobachtungen durch gehörige Reductions-Rechnung abgeleitet waren.
Freilich enthält, wie vorhin gezeigt, die ganze Anklage gar keinen Grund, dies zu
bezweifeln; freilich ist, auch ganz abgesehen von dieser Anklage, gar kein Grund zu
einem solchen Zweifel vorhanden, der in sich selbst schon darum ungereimt wäre,
weil es viel weniger Arbeit kostet, die wirklichen Beobachtungen zu reduciren, als
die Cometenörter aus Elementen zu berechnen. Allein diesmal giebt wirklich die
Beobachtung selbst den evidentesten Beweis des ächten Ursprunges der von
\textsc{Pasquich} bekannt gemachten Cometenpositionen. Für die meisten Leser wird
schon die Geringfügigkeit der oben gefundenen Abweichungen ein solcher Beweis
sein; allein ein viel \SourceLeaf{642}\enlargethispage{1pt} stärkerer liegt noch in ihrem \emph{regelmässigen Gange}. Unsre Reductionselemente können bedeutend
verschieden sein von den wahren; das wird einigen, obwohl immer nur einen kleinen, Einfluss auf die
Cometenpositionen haben müssen. Die Vergleichung der Resultate der Rechnung nach zwei verschiedenen
Systemen von Reductionselementen, auf eine und dieselbe Reihe von Beobachtungen angewandt, wird
also kleine Unterschiede zeigen, aber Unterschiede, die nothwendig einem regelmässigen Gange
folgen. Die Vergleichung von Positionen hingegen, die aus Elementen berechnet wären, mit solchen,
die aus wirklichen Beobachtungen abgeleitet sind, würde Unterschiede geben, die in Rücksicht des
Absoluten von den Fehlern der Elemente, und in Rücksicht des Relativen von den unordentlichen
Beobachtungsfehlern (verbunden mit den etwa absichtlich und willkührlich angebrachten kleinen
Abänderungen, wenn man solche annehmen wollte) die sichere Spur zeigen müssten, und es wäre mehr
als ein Wunder, wenn jemand, der so unvernünftig wäre, nach Elementen Positionen zu erdichten,
während er gute leicht zu reducirende Beobachtungen vor sich hat, bei einem solchen thörichten
Betruge solch eine Quinterne aus dem Glückstopfe zöge, dass er haarscharf dasselbe träfe, was ihm
die Reduction seiner Beobachtungen gegeben haben würde. Wie gross übrigens die Regelmässigkeit in
dem Gange der obigen Unterschiede ist, wird man am besten übersehen, wenn man die
Reductionselemente nicht aus den Beobachtungen von $\gamma$ Pegasi sondern aus denen des Cometen
selbst, verglichen mit den bekannt gemachten Positionen, ableitet, und zwar nur aus zwei
Beobachtungen des Cometen, um dann nachzusehen, wie die dritte damit harmonirt. Man findet auf diese Weise\HistoricalNote{36}:

\begin{center}\footnotesize\setlength{\tabcolsep}{3pt}
\begin{tabular}{@{}l r@{}}
Abstand des Pols des Instruments vom wahren Weltpole & $2'\;32''41$\\
Stundenwinkel des erstern . . . . . . . . . . & $170^\circ\;25\;17''$\\
Correction des Index für den Stundenwinkel . . . & $-3\;23.47$\\
Correction des Index für die Declination . . . . . & $-1\;34.38$
\end{tabular}\end{center}

Die auf ähnliche Art wie oben ausgeführte Reduction der Beobachtungen des Cometen gibt hienach

\begin{center}\footnotesize\setlength{\tabcolsep}{3pt}\begin{tabular}{r | r r | r | r r | r}
 & \multicolumn{2}{c|}{Stundenwinkel} & nach \textsc{Pasquich} & \multicolumn{2}{c|}{Declination} & nach \textsc{Pasquich} \\
\hline
Febr. 22 & $74^\circ\,44'$ & $8''99$ & $74^\circ\,44'\,9''$ & $14^\circ\,38'$ & $9''82$ & $14^\circ\,38'\,10''$ \\
26 & $81$ & $6\ \ 32.98$ & $81\ \ 6\ \ 33$ & $14$ & $23\ \ 13.36$ & $14\ \ 23\ \ 13$ \\
27 & $85$ & $10\ \ 10.03$ & $85\ \ 10\ \ 10$ & $14$ & $18\ \ 44.81$ & $14\ \ 18\ \ 45$ \\
\end{tabular}\end{center}

also wirklich \emph{vollkommene} Uebereinstimmung, da \textsc{Pasquich} keine Brüche von Secunden
angesetzt hat.

Dieses System von Reductionselementen stellt die Beobachtungen von $\gamma$ Pegasi,
wenn auch nicht ganz so nahe, wie das obige, doch immer noch nahe genug dar: die Unterschiede
werden

\begin{center}\footnotesize\setlength{\tabcolsep}{4pt}\begin{tabular}{r | r | r}
 & \shortstack{im Stunden-\\winkel} & \shortstack{in der\\Declination} \\
\hline
Febr. 20 & $+3''51$ & $-0''39$ \\
22 & $+12.47$ & $-10.32$ \\
22 & $+4.08$ & $-9.82$ \\
27 & $+10.81$ & $-11.42$ \\
27 & $-6.80$ & $-5.10$ \\
\end{tabular}\end{center}

Die von \textsc{Pasquich} wirklich angewandten Reductionselemente kennen wir nicht; sie können
bedeutend verschieden sein von den eben angeführten, da es zu misslich ist, aus den Correctionen auf \SourceLeaf{643} deren Elemente zurückzuschliessen, zumal da wir die Correctionen wegen des Instruments nicht von der Refraction getrennt aus den bekannt gemachten Datis erhalten können, und \textsc{Pasquich} vermuthlich die Refractionen nach andern Tafeln und vielleicht nach andern Methoden berechnet hat. Auch gehört die Frage, ob die Reductionselemente, welche \textsc{Pasquich} angewandt hat, die möglich genauesten gewesen sind, gar nicht zur Sache; oder vielmehr, diese Frage blos aufwerfen, heisst schon, \textsc{Pasquich} von der ihm gemachten Beschuldigung frei sprechen. Indessen erkläre ich gern, dass ich gar keinen Grund sehe, zu bezweifeln, dass \textsc{Pasquich} seine Reductions-Elemente mit gutem Vorbedacht aus wahrscheinlich viel zahlreichern und vielleicht ganz andern Beobachtungen abgeleitet habe, da eben aus obiger Berechnung selbst vollkommen erhellt, dass \textsc{Pasquich} wirklich seine Beobachtungen als absolute und nicht als Differential-Beobachtungen reducirt hat.

Ich habe bisher die Sache blos in wissenschaftlicher Beziehung betrachtet, wie könnte man aber unterlassen, sie auch aus dem Gesichtspunkte der Ehre und Rechtlichkeit anzusehen und ganz die gerechte Indignation zu theilen, die der edle \textsc{Olbers} so treffend ausgesprochen hat!

\noindent Göttingen 1824 März 3. \hfill C.\,F.\ \textsc{Gauss}.

[Handschriftliche Bemerkung: \textsc{Daniel Kmeth}, geb. zu Bries (Brezno Bamja) in Ungarn 1783 Jan. 15, starb in Kaschau 1825 Juni 20.]
